\chapter{Resolution of Cosmological Anomalies: Dark Energy, Dark Matter, and the Hubble Tension}
\label{chap:cosmological_anomalies}

\section{The Hubble Artifact and the Viscous Horizon}
When Edwin Hubble discovered the linear relationship between redshift and distance, standard physics mathematically misinterpreted this optical friction as a recessional velocity ($v = H_0 d$). Under TFT, the Hubble constant ($H_0$) is not a rate of spatial expansion; it is strictly the kinematic dissipation rate of the vacuum ($\Gamma_D \equiv H_0$).

This exposes a profound misconception in standard cosmology regarding both the ``Age of the Universe'' and the origin of the CMB. Standard physics calculates the age of the universe as $\approx 13.8$ billion years by taking the inverse of the Hubble constant ($1/H_0$). When telescopes observe a uniform wall of microwave radiation at this exact boundary, standard cosmology assumes it is the ``Surface of Last Scattering''---the primordial flash of the Big Bang.

In TFT, because $H_0$ is a coefficient of friction, its inverse ($\tau_{\text{visc}} = 1/H_0$) is simply the Photon Decay Constant.

It takes approximately 13.8 billion years for the hydrodynamic drag of the lattice to completely drain a photon's coherent wave energy.

At this exact optical distance, the photon ``dies.'' Its transverse shear wave completely loses topological coherence and dissolves into the incoherent, longitudinal phonon vibration of the lattice.

Therefore, 13.8 billion years is not the age of the universe, and the distant microwave wall is not the beginning of time. It is the Viscous Horizon---the strict optical boundary where redshifted starlight fully thermalizes into the Cosmic Microwave Background. We cannot see further back in time because the lattice is not perfectly transparent; standard physics is simply measuring the exact distance at which light physically drowns in the thermal noise floor of the vacuum.

\section{The ``Accelerating'' Expansion Illusion (The Logarithmic Law)}
In 1998, observations of distant Type Ia supernovae revealed that they were significantly dimmer and more redshifted than the linear Hubble law predicted \cite{Riess_1998}. To explain this anomaly within an expanding metric, standard cosmology invented ``Dark Energy'' to force the expansion of the universe to accelerate.

Tension Field Theory eliminates the need for Dark Energy by analyzing the supernova data through standard visco-elastic fluid dynamics. In any physical medium, the energy loss of a wave subjected to constant friction follows a compounding exponential decay function:
\begin{equation}
    E(x) = E_0 e^{-\frac{\Gamma_D x}{c_E}}
\end{equation}

Because observational redshift $(1+z)$ is the direct ratio of emitted energy to received energy, the true coordinate distance ($d$) of a wave subjected to this constant friction algebraically rearranges into a logarithmic attenuation curve. Substituting the observational kinematic drag constant ($\Gamma_D \equiv H_0$), the astronomical coordinate distance is:
\begin{equation}
    d = \frac{c_E}{H_0} \ln(1+z)
\end{equation}

However, astronomers do not measure raw coordinate distance; they measure Luminosity Distance ($d_L$). As derived in \hyperref[chap:cosmic_redshift]{Chapter \ref*{chap:cosmic_redshift}}, the topological conservation of the wave packet requires the signal to undergo temporal broadening, diluting the observed flux by a geometric factor of $(1+z)$. Therefore, the observed luminosity distance is the product of the logarithmic drag and the temporal dilution:
\begin{equation}
    d_L = \frac{c_E}{H_0} (1+z)\ln(1+z)
\end{equation}

When we analyze the Taylor series expansion of this exact visco-elastic luminosity distance, it exposes the fundamental mathematical error of standard cosmology. Expanding $(1+z)\ln(1+z)$ yields $z + \frac{z^2}{2}$. Over relatively local cosmic distances, the linear term ($z$) dominates, which is why the linear Hubble Law perfectly describes nearby galaxies. However, over extreme optical distances, the non-linear compounding of the lattice drag becomes dominant, revealing a positive second-order term ($+\frac{z^2}{2}$).

When standard cosmology attempts to force this natural drag curve into a purely kinematic Doppler-expansion model, the math misinterprets this exact $+\frac{z^2}{2}$ geometric term as a positive kinematic acceleration. The universe is not accelerating outward; the mathematical signature of Dark Energy is simply the standard Taylor expansion of visco-elastic boundary friction compounding over extreme distances (See \hyperref[app:hubble_diagram]{Appendix \ref*{app:hubble_diagram}} for the full algebraic derivation of the TFT Distance Modulus).

\section{The Hubble Tension: Acoustic Drag and the Standard Ruler}
The most pressing discrepancy in modern cosmology is the ``Hubble Tension''---an irreconcilable statistical divergence between the cosmic expansion rate measured locally ($H_0 \approx 73.0 \text{ km/s/Mpc}$) \cite{SH0ES_2022} and the rate inferred from the Cosmic Microwave Background ($H_0 \approx 67.4 \text{ km/s/Mpc}$) \cite{Planck_2018}. Standard cosmology attempts to resolve this discrepancy by postulating novel, unobserved physics (e.g., ``Early Dark Energy'').

TFT resolves this tension strictly through classical thermodynamics and solid-state Euclidean geometry. Standard cosmology derives the early-universe expansion rate by measuring the angular radius ($\theta_*$) of acoustic peaks and applying geometric triangulation against a theoretical ``Standard Ruler''---the Sound Horizon ($r_s$). However, standard models calculate $r_s$ by assuming the primordial plasma was a frictionless, inviscid fluid expanding within a relativistic metric. By integrating this expanding metric (yielding a comoving distance to the CMB of $D_M \approx 14,000 \text{ Mpc}$), standard cosmology mathematically inflates the theoretical physical radius of the sound horizon to $r_s \approx 145 \text{ Mpc}$.

\subsection{Viscous Acoustic Damping in the Active Vacuum}
In the TFT framework, the primordial plasma was not expanding through an empty geometric void; it was hydrodynamically coupled to the visco-elastic Spacetime Gravitational Floor (SGF).

\begin{itemize}
    \item \textbf{The Waves:} While transverse optical waves propagate through the solid FCC nodes, the Baryon Acoustic Oscillations (BAO) were longitudinal pressure waves propagating through the baryonic plasma and the interstitial condensate.
    \item \textbf{The Drag:} The propagation of these longitudinal acoustic waves was directly opposed by the inherent Lattice Impedance Acceleration ($a_D$). As formally derived in \hyperref[app:cosmic_impedance]{Appendix \ref*{app:cosmic_impedance}} ($a_D = c_E \Gamma_D$), this impedance is the fundamental material resistance of the SGF---the exact same reactive mechanism responsible for flat galactic rotation curves (the ``Dark Matter'' anomaly).
    \item \textbf{The Result:} This impedance acts as a continuous macroscopic boundary friction on the oscillating wave, restricting the physical distance the sound waves could travel before lattice annealing froze them into place.
\end{itemize}

Consequently, the actual physical density structures are significantly more compact than frictionless standard models predict.

\subsection{The Euclidean Acoustic Resolution}
Because TFT Vacuum operates as a static, visco-elastic lattice, the physical geometry of these density structures is absolute and immune to relativistic metric distortions. The physical optical limit of the cosmos is determined directly by the uncalibrated kinematic baseline observed locally by SH0ES ($H_0 = 73.0 \text{ km/s/Mpc}$). As established in \hyperref[app:vacuum_damping]{Appendix \ref*{app:vacuum_damping}}, this empirical measurement is not a metric expansion velocity, but the direct macroscopic manifestation of the theoretical Vacuum Damping Rate ($\Gamma_D \equiv H_0$). This temporal lattice friction sets the absolute optical boundary of the active vacuum:
\begin{equation}
    R_U = \frac{c_E}{H_0} \approx 13.39 \text{ Gly} \approx 4105 \text{ Mpc}
\end{equation}

We apply standard flat-space trigonometry to this absolute optical horizon to derive the true physical radius ($r_s$) of the primordial acoustic wave, utilizing the exact angular radius measured by the Planck satellite ($\theta_* \approx 0.01041 \text{ rad}$):
\begin{equation}
    r_{s(\text{TFT})} = R_U \cdot \theta_*
\end{equation}
\begin{equation}
    r_{s(\text{TFT})} = 4105 \text{ Mpc} \cdot 0.01041 \text{ rad} \approx 42.7 \text{ Mpc}
\end{equation}

In a visco-elastic Euclidean universe, the true acoustic ruler is exactly $42.7 \text{ Mpc}$. Standard cosmology calculates a massive $\approx 145 \text{ Mpc}$ structure solely because it mathematically forces this compact, damped object into an expanding relativistic metric.

\subsection{Recalibrating the Kinematic Baseline}
Applying the Euclidean standard ruler to the observed Planck geometry entirely eliminates the cosmological discrepancy. By correcting the frictionless assumption and metric expansion, the sequence of cosmological derivation aligns perfectly with local observations:

\begin{itemize}
    \item \textbf{Standard Model (Expanding Metric):} Frictionless Plasma $\to$ Inflated Comoving Distance ($D_M \approx 14,000 \text{ Mpc}$) $\to$ Inflated Ruler ($r_s \approx 145 \text{ Mpc}$) $\to$ Inferred Rate ($H_0 \approx 67.4 \text{ km/s/Mpc}$).
    \item \textbf{TFT (Static Viscous Lattice):} Damped Plasma $\to$ Absolute Optical Horizon ($R_U \approx 4105 \text{ Mpc}$) $\to$ True Euclidean Ruler ($r_s \approx 42.7 \text{ Mpc}$) $\to$ True Decay Rate ($H_0 \approx 73.0 \text{ km/s/Mpc}$).
\end{itemize}

The Hubble Tension is not a cosmological crisis requiring novel particle physics. It is the exact, quantifiable signature of primordial acoustic drag. The discrepancy is strictly a calibration error resulting from standard cosmology measuring a static, visco-elastic universe with an expanding, frictionless ruler.

\begin{figure}[H]
    \centering
    \includegraphics[width=0.9\textwidth]{Fig21_PartIII_CMB.pdf}
    \caption{\textbf{The Macroscopic Thermodynamic Projection \& Viscous Acoustic Drag.} A Euclidean geometric representation of Cosmic Microwave Background (CMB) anisotropies within a static, visco-elastic lattice.
    \newline \textbf{The Angular Projection:} The observed angular radius ($\theta_* \approx 0.01041 \text{ rad}$) projects from the Local Detector onto the Absolute Optical Horizon ($R_U \approx 13.39 \text{ Gly}$).
    \newline \textbf{The Horizon Limit:} The horizon distance is strictly determined by the Vacuum Damping Rate ($\Gamma_D \equiv H_0 = 73.0 \text{ km/s/Mpc}$).
    \newline \textbf{The Standard Ruler Resolution:} This static configuration yields a True Acoustic Ruler of $r_s \approx 42.7 \text{ Mpc}$, identifying the ``Hubble Tension'' as a calibration artifact resulting from the Standard Model's utilization of an inflated, frictionless ruler ($\approx 145 \text{ Mpc}$) within a mathematically expanding metric.
    \newline \textbf{Thermal Anisotropies:} The $\pm 30\mu\text{K}$ anisotropies are interpreted as Thermodynamic Impedance Shadows---thermal deficits cast by foreground density structures upon the local phonon screen of the vacuum.}
    \label{fig:cmb_acoustic_drag}
\end{figure}

\section{The Thermodynamic Impedance Signature and the $1.04^\circ$ Anomaly}
Standard cosmology assumes the characteristic $100\theta_* \approx 1.041$ angular scale measured by the Planck satellite \cite{Planck_2018} represents direct images of primordial acoustic waves frozen in a distant CMB plasma. However, because TFT establishes the CMB as the omnipresent longitudinal phonon heat of the local vacuum, these angular variations require a visco-elastic reinterpretation.

The fossilized acoustic rings are physically existing density structures located exactly at the absolute optical horizon ($R_U$). As established in Part I (Section 3.1), for a transverse wave to survive the $R_U$ distance and reach the local detector before dissolving into the thermal noise floor, it must have originated with the absolute maximum kinematic energy of the lattice ($E_{\text{max}}$). Normal optical energy emitted by standard baryonic processes at this extreme distance would dissipate long before reaching the observer.

Therefore, the satellite is not seeing normal thermal emission. It is detecting the degraded $E_{\text{max}}$ wave knots originally generated by the extreme localized shear of the omnipresent Big Rupture (Catastrophic Decompression). Because the Rupture occurred simultaneously across the entire $10^{35} \text{ m}$ cosmic reservoir, the satellite is specifically capturing the shell of $E_{\text{max}}$ wavefronts that originated at a distance of $R_U$---the exact distance required for vacuum drag to degrade $E_{\text{max}}$ to $2.725 \text{ K}$ upon arrival.

When this $E_{\text{max}}$ energy propagates from the baseline active vacuum at $R_U$, it decays to exactly $2.725 \text{ K}$. However, when this $E_{\text{max}}$ energy originates from within one of the dense, fossilized baryonic rings, the wave must overcome the elevated visco-elastic friction (the local $\alpha_G$ strain) of that structure. Because it pays this initial drag penalty, it begins its cosmic journey with fractionally less kinematic capacity. Consequently, it entirely dissipates into the phonon noise floor slightly colder than the baseline. The $\pm 30\mu\text{K}$ spots mapped by CMB satellites are therefore localized thermal deficits---Thermodynamic Impedance Signatures---perfectly reflecting the local gravitational strain of the $42.7 \text{ Mpc}$ acoustic rings without violating the bulk macroscopic damping rate ($\Gamma_D$).

\section{Kinematic Unification: Lattice Impedance ($a_D$) and the Dark Sector}
A primary theoretical synthesis of the Tension Field Theory framework is the kinematic unification of the Dark Sector.

In Part I, we derived the structural energy budget of the cosmos strictly from Face-Centered Cubic (FCC) geometry. We identified Dark Energy ($\beta_\Lambda$) as the effective structural packing limit of the solid lattice nodes ($\approx 69\%$), and Dark Matter ($\beta_{IC}$) as the interstitial condensate filling the voids ($\approx 26\%$). While these represent two distinct structural phases of the vacuum, the observational phenomena attributed to them in standard cosmology are both kinematic illusions produced by the mechanical interaction at their boundary layer.

Standard cosmology treats the ``accelerating'' cosmic redshift (Dark Energy) and flat galactic rotation curves (Dark Matter) as two entirely unrelated exotic phenomena. In TFT, they are the exact same boundary friction manifesting at different scales.

As established earlier, the Lattice Damping Rate ($\Gamma_D \equiv H_0$) represents the viscous energy loss of an optical wave caused by the solid lattice nodes displacing the interstitial condensate. Multiplying this visco-elastic phase lag by the elastic wave speed of the vacuum ($c_E$) yields the fundamental Lattice Impedance Acceleration ($a_D$), as formally derived in \hyperref[app:cosmic_impedance]{Appendix \ref*{app:cosmic_impedance}}:
\begin{equation}
    a_D = c_E \cdot \Gamma_D \approx 6.5 \times 10^{-10} \text{ m/s}^2
\end{equation}

In Part II, we established that the anomalous flat rotation curves of galaxies are strictly the inductive reactance of this exact same vacuum impedance acting upon baryonic mass. Observational astrophysics (via Modified Newtonian Dynamics, or MOND) has long measured this galactic minimum acceleration limit as $a_0 \approx 1.1 \times 10^{-10} \text{ m/s}^2$ \cite{Milgrom_1983}.

These two kinematic anomalies are mechanically identical, linked precisely by the spherical geometric projection factor of a galactic halo ($2\pi$):
\begin{equation}
    a_0 = \frac{a_D}{2\pi}
\end{equation}

\textbf{Conclusion:} The universe does not require an exotic anti-gravity field to push it apart, nor invisible particles to hold its galaxies together. The observable effects of the Dark Sector are strictly the result of boundary friction between the solid lattice nodes and the interstitial condensate. This energy dissipation by FCC lattice damping rate ($\Gamma_D$) acting upon a transverse optical wave creates the illusion of accelerating space; that exact same friction manifesting as inductive reactance ($a_D$) against the hydrodynamic inertia of orbiting stars creates flat rotation curves.